% =====================================================================
% LAMPIRAN F: LANDASAN TEORI PEMROSESAN MULTIMODAL DAN REDUKSI DIMENSI
% =====================================================================
\chapter{Landasan Teori Pemrosesan Multimodal dan Reduksi Dimensi}
\label{app:teori-multimodal}

Lampiran ini menampung tiga bagian landasan teori yang sebelumnya berada di Bab~\ref{ch:tinjauan} namun tidak berkontribusi langsung sebagai keputusan metodologi pada Bab~\ref{ch:metode}. Pemindahan dilakukan untuk menjaga pembahasan Bab~\ref{ch:tinjauan} tetap fokus pada teori yang dipakai pada pipeline akhir, sementara latar teori yang lebih umum tetap dapat ditinjau di sini.

\section*{F.1 Landasan Teori Pemrosesan Data Gambar}
\label{app:imgproc-detail}

Bagian ini menguraikan kerangka teoretis pra-proses citra pada \textit{pipeline} klasifikasi multimodal. Pada penelitian ini seluruh pra-proses citra didelegasikan kepada \textit{AutoImageProcessor} dari pustaka HuggingFace, khususnya kelas \textit{DINOv3ViTImageProcessorFast}, yang dimuat sekali pada saat inisialisasi server FastAPI dan mengikuti spesifikasi resmi yang dipublikasikan bersama bobot DINOv3. Konfigurasi pra-proses tersimpan pada berkas \texttt{preprocessor\_config.json} dan menjadi sumber kebenaran (\textit{source of truth}) untuk seluruh tahap inferensi. Subbagian berikut memaparkan komponen-komponen yang aktif pada konfigurasi tersebut beserta latar teoretisnya.

\subsection*{F.1.1 Citra Digital sebagai Tensor}

Citra digital adalah representasi visual objek yang diubah ke dalam format numerik melalui perangkat akuisisi data seperti kamera. Angka tersebut mencerminkan intensitas cahaya pada setiap elemen gambar (\textit{pixel}), yang memungkinkan analisis dan manipulasi melalui teknik pengolahan citra digital \citep{article}. Pada \textit{deep learning}, satu citra berwarna direpresentasikan sebagai tensor berukuran $H \times W \times C$, dengan $H$ tinggi, $W$ lebar, dan $C$ adalah jumlah kanal warna (umumnya 3 untuk RGB). Nilai piksel umumnya berada pada rentang bulat $[0, 255]$ untuk berkas mentah \texttt{uint8}, atau pada rentang \textit{floating point} $[0, 1]$ atau $[-1, 1]$ setelah normalisasi.

\subsection*{F.1.2 \textit{Resize} dan Interpolasi}

Karena \textit{encoder} citra menerima masukan dengan ukuran tetap, setiap citra harus dilakukan \textit{resize} ke resolusi target. Tiga algoritma interpolasi yang paling umum adalah:

\begin{itemize}
    \item \textbf{Nearest Neighbor}: piksel keluaran mengambil nilai piksel terdekat dari masukan. Cepat namun menghasilkan tepi bergerigi.
    \item \textbf{Bilinear}: rata-rata berbobot dari empat piksel tetangga. Halus dan menjadi standar pada banyak \textit{pipeline}.
    \item \textbf{Bicubic}: rata-rata berbobot dari enam belas piksel tetangga dengan polinomial kubik. Lebih tajam namun lebih mahal.
\end{itemize}

\textit{AutoImageProcessor} pada DINOv3 melakukan \textit{resize} ke resolusi target $224 \times 224$ piksel mengikuti parameter \texttt{size} dan \texttt{interpolation} yang ditetapkan oleh kartu model resmi.

\subsection*{F.1.3 Aspek Rasio dan \textit{Square Padding}}

Pada konfigurasi \textit{AutoImageProcessor} DINOv3 berlaku \texttt{default\_to\_square: true} dengan \texttt{do\_center\_crop} dinonaktifkan, sehingga citra dipaksa berbentuk persegi melalui kombinasi \textit{resize} dan \textit{padding} pada sisi yang lebih pendek. Strategi ini berbeda dengan \textit{center crop} klasik yang memotong tepi citra setelah \textit{resize} sisi pendek; \textit{square padding} mempertahankan seluruh konten citra namun menambahkan piksel netral pada sisi yang kekurangan ukuran. Pada konteks laporan warga, pendekatan ini menguntungkan karena objek penting (lubang jalan, pohon tumbang, sampah) sering kali tersebar tidak terpusat sehingga \textit{cropping} tepi berisiko membuang informasi diskriminatif.

\subsection*{F.1.4 Normalisasi Per-Kanal}

Setelah \textit{resize}, citra dinormalisasi per kanal menggunakan rerata dan simpangan baku ImageNet:

\begin{equation}
    x'_c = \frac{x_c - \mu_c}{\sigma_c}, \quad
    \mu = (0{,}485,\, 0{,}456,\, 0{,}406), \;
    \sigma = (0{,}229,\, 0{,}224,\, 0{,}225),
    \label{eq:norm}
\end{equation}
dengan $c$ adalah indeks kanal RGB. Nilai $\mu$ dan $\sigma$ ini termuat secara eksplisit pada \texttt{preprocessor\_config.json} (kunci \texttt{image\_mean} dan \texttt{image\_std}) dan wajib digunakan ketika \textit{encoder} DINOv3 pra-latih, agar distribusi masukan pada saat inferensi konsisten dengan distribusi yang dipakai pada saat pra-pelatihan.

\subsection*{F.1.5 \textit{Layout} Tensor dan Konversi Warna}

Dua \textit{layout} tensor yang umum adalah NCHW (\textit{batch, channel, height, width}) yang dipakai PyTorch dan ONNX, serta NHWC yang dipakai TensorFlow. Selain itu, beberapa pustaka (contohnya OpenCV) menggunakan urutan kanal BGR sementara model \textit{deep learning} biasanya mengasumsikan RGB. \textit{AutoImageProcessor} pada penelitian ini diatur dengan \texttt{data\_format: channels\_first} sehingga keluaran tensor langsung berformat NCHW. Konversi kanal ke RGB dilakukan secara eksplisit melalui pemanggilan \texttt{Image.open(...).convert("RGB")} pada lapisan \textit{encoder} citra sehingga sumber gambar berformat lain (BGR, RGBA, atau \textit{grayscale}) tetap menghasilkan tensor masukan yang konsisten.

\subsection*{F.1.6 Kompresi JPEG dan EXIF}

Berkas JPEG biasanya sudah dikompresi dengan kerugian (\textit{lossy}) dan dilengkapi metadata EXIF, termasuk orientasi sensor kamera. Pustaka PIL yang digunakan oleh \textit{AutoImageProcessor} membaca dan menerapkan tag orientasi EXIF secara otomatis melalui pemanggilan \texttt{Image.open()}, sehingga foto vertikal yang diambil dari kamera ponsel tidak dilewatkan ke model dalam orientasi miring. Kompresi JPEG yang berlebihan juga dapat menurunkan akurasi model, namun pada \textit{pipeline} produksi penelitian ini gambar yang diunggah pengguna disimpan langsung pada Supabase Storage sebagai berkas asli sehingga tidak terjadi kompresi tambahan setelah pengambilan citra.

\subsection*{F.1.7 Catatan tentang Augmentasi Data}

Teknik augmentasi data seperti \textit{horizontal flip}, \textit{color jitter}, \textit{random erasing}, RandAugment \citep{cubukRandAugment2020}, dan MixUp \citep{zhangMixup2018} merupakan komponen standar pada pelatihan \textit{encoder} citra dari nol. Pada penelitian ini teknik tersebut tidak diterapkan karena \textit{encoder} DINOv3 dibekukan (\textit{frozen}) selama seluruh tahap eksperimen dan hanya dipakai pada modus inferensi (\texttt{is\_training=False}). \textit{Embedding} dihitung sekali per sampel dan disimpan pada berkas \textit{cache} sehingga ablasi pada \textit{classifier head} dapat dijalankan tanpa perlu menjalankan kembali \textit{encoder}. Augmentasi data tetap dibahas di sini sebagai latar teoretis karena merupakan teknik penting yang relevan ketika DINOv3 di masa depan akan dilatih ulang penuh atau di-\textit{fine-tune} pada domain spesifik laporan warga.

\section*{F.2 Landasan Teori Pemrosesan Data Teks}
\label{app:textproc-detail}

\subsection*{F.2.1 Pembersihan dan Normalisasi}

Sebelum dilewatkan ke \textit{encoder} teks, sebuah dokumen biasanya melalui beberapa langkah pembersihan, yaitu:

\begin{itemize}
    \item Normalisasi unicode ke bentuk NFC (\textit{Canonical Composition}) agar karakter ber-aksen direpresentasikan secara konsisten.
    \item Konversi huruf ke \textit{lowercase} jika model toleran terhadap kasus, atau dipertahankan apabila model membedakan kasus.
    \item Penghapusan karakter kontrol dan \textit{whitespace} berlebih (\textit{whitespace collapsing}).
    \item Substitusi emoji yang sangat jarang dengan \textit{placeholder}, untuk menghindari pembengkakan kosakata.
    \item Pemotongan teks menjadi panjang maksimum, biasanya 128 atau 512 token, sesuai batas \textit{encoder}.
\end{itemize}

\subsection*{F.2.2 Tokenisasi}

Tokenisasi adalah proses memecah teks menjadi unit diskrit yang dapat dipetakan ke indeks. Tiga keluarga utama tokenizer adalah:

\begin{enumerate}
    \item \textbf{Word-level}: setiap kata adalah satu token. Sederhana namun rentan terhadap kosakata yang membengkak dan kata di luar kosakata (\textit{out-of-vocabulary}).
    \item \textbf{Character-level}: setiap karakter adalah satu token. Tahan terhadap \textit{out-of-vocabulary}, namun barisan menjadi sangat panjang.
    \item \textbf{Subword}: kompromi antara dua pendekatan di atas, yaitu memecah kata menjadi unit subkata yang sering muncul. Ini adalah pilihan dominan pada \textit{Transformer} modern.
\end{enumerate}

Empat algoritma subword yang paling umum:

\begin{itemize}
    \item \textbf{Byte Pair Encoding} (BPE) \citep{sennrichBPE2016}: mulai dari karakter sebagai unit, secara iteratif menggabungkan pasangan unit yang paling sering bersamaan menjadi satu unit baru hingga ukuran kosakata target tercapai.
    \item \textbf{WordPiece}: varian BPE yang dipakai oleh BERT \citep{devlinBERT2019}, dengan kriteria \textit{merge} berbasis \textit{likelihood} alih-alih frekuensi.
    \item \textbf{SentencePiece} \citep{kudoSentencePiece2018}: \textit{wrapper} BPE atau Unigram yang bekerja pada level byte sehingga tidak memerlukan pra-proses spesifik bahasa, cocok untuk tokenizer multibahasa.
    \item \textbf{Unigram Language Model}: memodelkan distribusi token dengan probabilitas independen, kemudian memilih segmentasi yang memaksimalkan \textit{likelihood}.
\end{itemize}

Tabel berikut merangkum perbandingan singkat keempatnya.

\begin{table}[htbp]
\centering
\caption{Perbandingan algoritma tokenizer subword.}
\label{tab:tokenizer}
\begin{tabular}{p{2.6cm} p{4cm} p{4cm} p{3cm}}
\toprule
\textbf{Algoritma} & \textbf{Kriteria \textit{merge}} & \textbf{Implementasi populer} & \textbf{Catatan} \\
\midrule
BPE & Frekuensi pasangan & GPT-2, RoBERTa & Sederhana, deterministik \\
WordPiece & \textit{Likelihood} pasangan & BERT, mBERT, IndoBERT & Memakai prefiks \texttt{\#\#} \\
SentencePiece & BPE atau Unigram di level byte & T5, mT5, XLM-R, E5 & Tidak butuh pra-tokenisasi spasi \\
Unigram LM & Probabilitas token & SentencePiece (mode Unigram) & Sampling segmentasi mungkin \\
\bottomrule
\end{tabular}
\end{table}

\subsection*{F.2.3 \textit{Embedding}: \textit{Lookup Matrix} dan Dimensi}

Setelah token diubah menjadi indeks bilangan bulat, indeks tersebut dipetakan ke vektor berdimensi $d$ melalui matriks \textit{lookup} berukuran $|V| \times d$, dengan $|V|$ ukuran kosakata. Vektor inilah yang menjadi masukan ke lapisan \textit{Transformer} berikutnya.

\section*{F.3 Reduksi Dimensi dan \textit{Principal Component Analysis}}
\label{app:pca-detail}

\subsection*{F.3.1 Motivasi}

Vektor representasi (\textit{embedding}) yang dihasilkan oleh \textit{encoder} modern umumnya berdimensi sangat tinggi, contohnya 768, 1024, atau bahkan 4096 untuk model besar. Pada banyak skenario, dimensi tinggi dapat menjadi masalah karena tiga alasan, yaitu (1) biaya memori dan komputasi pada model klasifikasi hilir, (2) \textit{curse of dimensionality} pada metode berbasis jarak (k-Nearest Neighbors, klastering jarak Euclidean), dan (3) sulitnya visualisasi serta interpretasi. Reduksi dimensi adalah keluarga teknik yang memetakan vektor berdimensi tinggi ke ruang berdimensi lebih rendah dengan menjaga sebanyak mungkin informasi yang relevan.

\subsection*{F.3.2 \textit{Principal Component Analysis} (PCA)}

PCA \citep{joliffePCA2002} adalah metode reduksi dimensi linear yang menemukan basis ortogonal baru sehingga proyeksi data ke basis tersebut memaksimalkan varians yang dipertahankan. Diberikan matriks data $\mathbf{X} \in \mathbb{R}^{n \times d}$ yang sudah dipusatkan (\textit{mean-centered}) sehingga setiap kolom memiliki rerata nol, matriks kovarians sampel dihitung sebagai:

\begin{equation}
    \mathbf{C} = \frac{1}{n-1} \mathbf{X}^\top \mathbf{X}.
    \label{eq:pca_cov}
\end{equation}

Vektor eigen $\mathbf{v}_1, \mathbf{v}_2, \dots, \mathbf{v}_d$ dari $\mathbf{C}$ yang diurutkan berdasarkan nilai eigen $\lambda_1 \geq \lambda_2 \geq \dots \geq \lambda_d$ disebut \textit{principal components}. Komponen pertama menjelaskan varians terbesar; komponen kedua menjelaskan varians terbesar berikutnya pada arah yang ortogonal terhadap komponen pertama, dan seterusnya. Reduksi ke $k$ dimensi dilakukan dengan memilih $k$ vektor eigen pertama yang dikumpulkan ke matriks $\mathbf{V}_k \in \mathbb{R}^{d \times k}$, kemudian memproyeksikan data:

\begin{equation}
    \mathbf{Z} = \mathbf{X} \mathbf{V}_k, \quad \mathbf{Z} \in \mathbb{R}^{n \times k}.
    \label{eq:pca_proj}
\end{equation}

Implementasi praktis sering menggunakan \textit{Singular Value Decomposition} (SVD) pada $\mathbf{X}$, yaitu $\mathbf{X} = \mathbf{U} \boldsymbol{\Sigma} \mathbf{V}^\top$, sehingga vektor eigen kovarians sama dengan kolom $\mathbf{V}$ dan nilai eigen sama dengan kuadrat nilai singular dibagi $n-1$. Pemilihan jumlah komponen $k$ umumnya dipandu oleh \textit{explained variance ratio}, yaitu rasio kumulatif $\sum_{i=1}^{k} \lambda_i / \sum_{i=1}^{d} \lambda_i$, dengan ambang umum 0,90 hingga 0,99.

\subsection*{F.3.3 Penggunaan dalam \textit{Pipeline Embedding}}

Pada \textit{pipeline} klasifikasi berbasis \textit{embedding}, PCA dapat dipakai untuk tiga tujuan:

\begin{itemize}
    \item \textbf{Pra-proses fitur}: mengurangi dimensi \textit{embedding} sebelum dilewatkan ke algoritma klasifikasi yang sensitif terhadap dimensi tinggi (contohnya regresi logistik dengan regularisasi atau k-NN).
    \item \textbf{Whitening}: PCA dengan normalisasi $1 / \sqrt{\lambda_i}$ pada setiap komponen menghasilkan ruang yang fitur-fiturnya tidak berkorelasi dan memiliki skala seragam.
    \item \textbf{Visualisasi}: memproyeksikan \textit{embedding} ke 2 atau 3 dimensi pertama untuk inspeksi klaster kualitatif, walaupun untuk visualisasi non-linear teknik t-SNE atau UMAP umumnya lebih informatif.
\end{itemize}

\subsection*{F.3.4 Keterbatasan dan Alternatif}

PCA bersifat linear, sehingga tidak dapat menangkap struktur non-linear yang sering muncul pada \textit{embedding} jaringan saraf dalam. Alternatif non-linear yang umum digunakan adalah \textit{Kernel PCA}, \textit{Autoencoder}, t-SNE, dan UMAP. Selain itu, PCA tidak memperhatikan label kelas; bila reduksi dimensi yang \textit{discriminative} terhadap kelas dibutuhkan, \textit{Linear Discriminant Analysis} (LDA) menjadi pilihan yang lebih tepat. Pada konteks pohon keputusan seperti CatBoost atau XGBoost, reduksi dimensi melalui PCA sering kali tidak memberi keuntungan karena pohon keputusan memilih dimensi yang relevan secara implisit melalui pemilihan \textit{split} pada setiap simpul.
